\documentclass[11pt]{article}
\usepackage[T1]{fontenc}
\usepackage[utf8]{inputenc}
\usepackage[english]{babel}
\usepackage{amsmath,amssymb,mathtools,array}
\usepackage{geometry}
\usepackage{microtype}
\geometry{a4paper,margin=3cm}
\setlength{\parindent}{1.25em}
\setlength{\parskip}{0pt}
\newcommand{\widebar}[1]{\overline{#1}}
\begin{document}

\begin{center}
{\Large\bfseries 2. On the Formation of the Form System of the Ternary\\ Biquadratic Form}\par
\vspace{1.5em}
Journal für die reine und angewandte Mathematik 134 (1908), pp. 23--90 and two tables
\end{center}

\vspace{3em}

\begin{center}
\textbf{Introduction.}
\end{center}

Work by Gordan, Maisano, and Pascal\footnote{A short extract from the Introduction and Chapter I appeared in the \emph{Sitzungsberichte der physikalisch-medizinischen Sozietät Erlangen 1907}, pp. 176--179.}\footnote{P. Gordan, on the complete form system of the ternary biquadratic form
\[
f=x_1^3x_2+x_2^3x_3+x_3^3x_1.
\]
(\emph{Math. Annalen}, vol. XVII (1880), pp. 217--233.)
G. Maisano, 1) \emph{Sistemi completi dei primi cinque gradi della forma ternaria biquadratica e degl' invarianti, covarianti e contravarianti di sesto grado}. (Giorn. di Battaglini XIX.)
2) \emph{Sui covarianti indipendenti di $6^\circ$ grado nei coefficienti della forma biquadratica ternaria}. (Rend. Circ. Mat. di Palermo I, 1887.)
E. Pascal, \emph{Contributo alla teoria della forma ternaria biquadratica e delle sue varie decomposizioni in fattori}. (Memoria premiata dalla R. Accademia delle scienze fisiche e matematiche di Napoli. 1905.)} deals with the form system of the ternary biquadratic form. Gordan sets up the complete form system, consisting of 54 formations, of the special automorphic form
\[
f=x_1^3x_2+x_2^3x_3+x_3^3x_1
\]
on the basis of principles similar to those that he had given for form systems in the binary domain.

In Maisano's work, for the general biquadratic form, the forms up to and including the fifth order\footnote{By ``order'' the dimension in the coefficients is to be understood, and by ``degree'' the dimension in the variables.} are set up, together with some invariants, covariants, and contravariants of higher order, according to the method used by Gordan in volume I of the \emph{Mathematische Annalen} for the ternary cubic form. Pascal, using Maisano's results, is chiefly concerned with the question of the decomposition of the biquadratic form into factors\footnote{In his second short note, Maisano attempts to prove a linear dependence of the three covariants of order 6 and degree 6. In contrast to this not quite complete proof, Pascal believes he has proved the linear independence of the three covariants of order 6, and likewise that of the three invariants of order 9. That, however, a linear relation does in fact exist between the three covariants on the one hand and the three invariants on the other is shown by formulas (13), § 11, and (22), § 17 note, of the present work, where these relations are given explicitly. According to a communication from Pascal, the contradiction is explained by a numerical error in the expression for his contravariant $p$ (called $\varrho$ here), p. 46 of his paper.}.

\emph{The aim of the following investigations is to set up the form system for the general ternary biquadratic form; in this work, namely, only the principal foundations are given, and a so-called ``relatively complete system''\footnote{For the terminology see § 3a.} is established.}

The work is closely connected with Gordan's work; however, the principles that were stated there only quite generally first had to be worked out in detail. With the aid of a theorem on the connection among the foldings and by introducing the ``form sequence'' (§ 1), the reduction theorems are sharply formulated and completed (§ 3), while the recursive establishment of special series expansions (§ 2) gives the computational means for actually carrying out the reductions.

The basic idea for forming systems of ternary forms is the same as in the binary domain. Starting from a first relatively complete system -- the system of forms taken over from the binary case -- one passes, according to a definite law, to systems with ever higher modulus, until either the system of a modulus -- the modulus taken as ground form -- becomes finite and known, or until a modulus can be reduced to forms having invariants as factors. By transvection of the relatively complete system with the system of the modulus, in the first case the absolutely complete system arises, while in the second case the relatively complete and absolutely complete systems coincide. By Hilbert's general proof of the finiteness of form systems, this procedure must necessarily come to an end.

In our case the modulus $(abc)$ of the first relatively complete system (§ 4) is reduced to the moduli
\[
\Delta=(abc)^2a_x^2b_x^2c_x^2\quad\text{and}\quad \nu=(abu)^4
\]
(§ 5), and then the relatively complete system modulo $\nu$ is found (§ 6). These results are already given in Gordan's work for the general biquadratic form; our manner of deriving them, however, is more easily transferred to forms of higher degree.

As the sequence of moduli we now choose the forms (§ 7)
\[
\nu,\qquad \nu(\nu)=(\nu\nu_1x)^4=s_x^4,\qquad \nu(s)=(ss'u)^4,\ldots .
\]

In forming the relatively complete system modulo $s$, two quadratic forms occurring in the system, $u_\varrho^2$ and $t_x^2$, are adjoined as moduli (§§ 9 and 10), and accordingly the relatively complete system modulo $(s,\varrho,t)$ is formed. It is then shown (§ 17) that the modulus $(ss'u)^4$ is reducible to the moduli $\varrho$ and $t$. The next higher system thereby passes over into a relatively complete system modulo $(\varrho,t)$. Since, however, the simultaneous system of two quadratic forms is finite and known, and in the general case consists of 20 formations\footnote{Cf. Clebsch--Lindemann, \emph{Vorlesungen über Geometrie}. (Vol. I, p. 288 ff.) System of two cogredient forms. Gordan, ``Über Büschel von Kegelschnitten.'' (\emph{Math. Ann.} vol. XIX, p. 530.) System of two contragredient forms.}, the termination described above is reached with the establishment of the relatively complete system modulo $(\varrho,t)$. The transvection of this system with the system of $(\varrho,t)$ in order to form the absolutely complete system is reserved.

As the relatively complete system modulo $(\varrho,t)$ one finds 331 formations, which in the appended table are ordered according to their degree in the variables $x$ and $u$.

Finally we give an overview of the symbols introduced:
\begin{align*}
f&=a_x^4,\\
\theta&=\theta_x^4u_\vartheta^2=(abu)^2a_x^2b_x^2,\\
K&=K_x^6u_k^3=(a\theta u)u_\vartheta^2a_x^3\theta_x^3,\\
j&=u_j^6=(a\theta u)^4u_\vartheta^2,\\
\Delta&=\Delta_x^6=a_\vartheta^2a_x^2\theta_x^4=(abc)^2a_x^2b_x^2c_x^2,\\
N&=N_x^8u_\eta=(a\Delta u)a_x^3\Delta_x^5,\\
\nu&=u_\nu^4=(abu)^4,\\
H&=H_x^2u_\eta^4=(\nu\nu_1x)^2u_\nu^2u_{\nu_1}^2,\\
L&=L_x^3u_l^6=(\nu\eta x)H_x^2u_\nu^3u_\eta^3,\\
g&=g_x^6=(\nu\eta x)^4H_x^2,\\
\sigma&=u_\sigma^6=H_\nu^2u_\nu^2u_\eta^4,\\
s&=s_x^4=(\nu\nu_1x)^4,\\
Z&=Z_x^4u_\zeta^2=(ss'u)^2s_x^2s_x'^2,\\
\varrho&=u_\varrho^2=\theta_\nu^4u_\vartheta^2,\\
t&=t_x^2=a_\eta^4H_x^2,\\
i&=a_\nu^4,\\
J&=s_\nu^4.
\end{align*}

\begin{center}
{\Large\bfseries Chapter I. \quad General Theorems on Ternary Forms.}
\end{center}

\begin{center}
\textbf{§ 1. \quad The folding process. \quad Form sequences.}
\end{center}

The fundamental process for producing forms is the folding process, which in the ternary domain can be defined as follows. Let there be given a symbolic product
\[
s_x^m t_x^n u_\sigma^\mu u_\tau^\nu .
\]
If one replaces the pairs of factors
\[
s_xt_x\qquad u_\sigma u_\tau\qquad s_xu_\sigma\ \text{or}\ s_xu_\tau
\qquad t_xu_\tau\ \text{or}\ t_xu_\sigma
\]
respectively by
\[
(stu)\qquad (\sigma\tau x)\qquad s_\sigma\ \text{or}\ s_\tau
\qquad t_\tau\ \text{or}\ t_\sigma,
\]
folding
\[
\text{I}\qquad\text{II}\qquad\text{III}\qquad\text{IV},
\]
then the resulting forms have arisen from the original one by folding\footnote{Gordan, \emph{Math. Annalen} vol. XVII, p. 219.}.

Here the expression $s_x^m t_x^n u_\sigma^\mu u_\tau^\nu$ can be regarded either as an actual product of two forms $S\cdot T$, or as a single form with several rows of symbols:
\[
s_x^m t_x^n u_\sigma^\mu u_\tau^\nu=A_x^{m+n}u_x^{\mu+\nu}.
\]
In the first case we speak of the folding of the form $S$ with $T$, in the second case of the ``folding of the form $A$ in itself.'' For brevity we always assume in what follows, as indeed is the case for all forms later to be considered, that
\[
\tag{a}
s_\sigma^\lambda=0,\qquad t_\tau^\lambda=0
\]
for any $\lambda$ and $x$ distinct from 0, and in conjunction with any other foldings. This says that the form $s_x^m t_y^n u_\sigma^\mu u_\tau^\nu$ is a so-called ``normal form,'' that is, it vanishes under application of the $\Omega$-process, both with respect to the variables $x$ and $u$ and with respect to the variables $y$ and $v$.

Thus condition (a) represents no restriction, but can always be achieved\footnote{Cf. Gordan, \emph{Math. Annalen} vol. V, p. 104.}.

For the connection among the individual foldings the following theorem holds:

\emph{Theorem I. Foldings I and II are fundamental foldings from which, independently of the order of composition, foldings III and IV can be composed. In other words: in order to form all forms arising by folding from a given expression, one has to apply only foldings I and II.}

Proof: By the identity theorem, respectively the product theorem for matrices, taking (a) into account, one obtains
\[
(\widehat{st}\sigma x)=-t_\sigma,\qquad
(\widehat{\sigma\tau}su)=-s_\tau,
\]
\[
(\widehat{st}\tau x)=s_\tau,\qquad
(\widehat{\sigma\tau}tu)=t_\sigma,
\]
\[
(stu)(\sigma\tau x)=s_\tau+t_\sigma-u_x\,s_\tau t_\sigma .
\]
Thus by composing foldings I and II one obtains foldings III and IV, and this holds both when applying folding II to folding I, i.e. to the factors $(stu)u_\sigma u_\tau$, and when applying folding I to folding II, to the factors $(\sigma\tau x)s_xt_x$.

By a \emph{form sequence}\footnote{Cf. Clebsch, ``Über eine Fundamentalaufgabe der Invariantentheorie'' (Abh. der Gött. Ges. d. Wiss. vol. XVII): § 17 and end of § 18. The ``form sequence'' differs from the ``properly reduced equivalent system'' introduced there by the principle of ordering according to higher forms.} we mean an initial form together with all forms that have arisen from it by folding in itself, and we denote the form sequence by the initial form. A ``higher form'' is to mean here any form with more folding in itself, while among the equally entitled foldings $s_\tau$ and $t_\sigma$ one has to be normalized as higher. According to the law of formation of the form sequence it follows that:

1) Each form of the form sequence is linear in the symbols of the initial form.

2) The form sequence is uniquely determined for initial forms with two rows each of contragredient symbols; for initial forms with more rows of symbols it is to be uniquely normalized by distinguishing special foldings among those connected by the identity theorem.

By Theorem I, a form sequence $s_x^m t_x^n u_\sigma^\mu u_\tau^\nu$ ($m>n;\ \mu>\nu$) can be arranged according to the following rectangular scheme:
\begingroup\small
\[
\begin{array}{c|c|c|c}
s_x^m t_x^n u_\sigma^\mu u_\tau^\nu
& (stu)
& (stu)^2
& \cdots\\[0.3em]
(\sigma\tau x)
& t_\sigma,\ s_\tau
& t_\sigma(stu),\ s_\tau(stu)
& \cdots\\[0.3em]
(\sigma\tau x)^2
& t_\sigma(\sigma\tau x),\ s_\tau(\sigma\tau x)
& t_\sigma^2,\ t_\sigma s_\tau,\ s_\tau^2
& \cdots\\[0.3em]
\vdots&\vdots&\vdots&\ddots\\[0.3em]
(\sigma\tau x)^\nu
& t_\sigma(\sigma\tau x)^{\nu-1},\ s_\tau(\sigma\tau x)^{\nu-1}
& t_\sigma^2(\sigma\tau x)^{\nu-2},\ t_\sigma s_\tau(\sigma\tau x)^{\nu-2},\ s_\tau^2(\sigma\tau x)^{\nu-2}
& \cdots\\[0.3em]
\vdots
& t_\sigma(\sigma\tau x)^\nu
& t_\sigma^2(\sigma\tau x)^{\nu-1},\ t_\sigma s_\tau(\sigma\tau x)^{\nu-1}
& \cdots
\end{array}
\]
\endgroup
and the correspondingly continued lower part\footnote{Here, and similarly in what follows, the lower part marked with an asterisk is to be set at the correspondingly marked place at the upper right of the scheme.}
\begingroup\small
\[
\begin{array}{c|c|c}
(stu)^n & \cdots & \cdots\\[0.3em]
t_\sigma(stu)^{n-1},\ s_\tau(stu)^{n-1}
& s_\tau(stu)^n & \cdots\\[0.3em]
t_\sigma^2(stu)^{n-2},\ t_\sigma s_\tau(stu)^{n-2},\ s_\tau^2(stu)^{n-2}
& t_\sigma s_\tau(stu)^{n-1},\ s_\tau^2(stu)^{n-1}
& s_\tau^2(stu)^n .
\end{array}
\]
\endgroup

Here one obtains, when one moves

1) one column to the right, all forms arising by folding I from the adjacent forms,

2) one row downward, all forms arising by folding II from the forms above,

and hence all forms arising by folding.

An arbitrary form of the total form sequence, $t_\sigma^\kappa s_\tau^\lambda(stu)^\mu$ or $t_\sigma^\kappa s_\tau^\lambda(\sigma\tau x)^\nu$, forms the beginning of a new form sequence, consisting of all those forms of the rectangular scheme with $t_\sigma^\kappa s_\tau^\lambda(stu)^\mu$ or $t_\sigma^\kappa s_\tau^\lambda(\sigma\tau x)^\nu$ as left upper corner, which have the factor $t_\sigma^\kappa s_\tau^\lambda$.

Supplementary remark. Theorem I has an exception in the case where the numbers $n$ and $\mu$ (respectively $m$ and $\nu$) become equal to zero, so that foldings I, II, and IV no longer exist, while folding III (respectively IV) still does. In this case we designate as the form sequence the diagonal member of the scheme:
\[
\begin{array}{ccccc}
s_x^m u_\tau^\nu &&&& t_x^n u_\sigma^\mu\\
& s_\tau && t_\sigma &\\
& s_\tau^2 && t_\sigma^2 &\\
&&\ddots&&\\
&&s_\tau^\nu && t_\sigma^\mu .
\end{array}
\]
In accordance with the absence of foldings I and II, the series expansion according to polars of the form sequence in this case always reduces to the initial member; however, the theorems on reducers remain valid.

\begin{center}
\textbf{§ 2. \quad Series expansions according to polars of the form sequence.}
\end{center}

For forms with $n$ variables two kinds of series expansions have been set up explicitly\footnote{Cf. Gordan, ``Über Kombinanten,'' §§ 2 and 5 (\emph{Math. Ann.} vol. V).}:

1) for forms with one row each of contragredient variables, $s_x^m u_\tau^\nu$, a series progressing by powers of $u_x$, whose coefficients are ``normal forms'';

2) for forms with two rows of cogredient variables, $s_x^m t_x^n$, an expansion according to polars of the elementary covariants (polars of the form sequence $s_x^m t_x^n$), which in the ternary case reads as follows\footnote{Cf. Study, \emph{Methoden zur Theorie der ternären Formen} (Teubner 1889) II. § 7. In distinction to the derivation there, ours gives the method for the rapid computational determination of the numerical coefficients $C_{pr\varrho}$. The Clebsch--Study expansion, upon specialization a), a'), would read
\[
(stu)^\mu u_\tau^{\nu-\kappa}v_\tau^\kappa s_x^m t_x^{\,n-\lambda}(tuv)^\lambda
=\sum_{p,r,\varrho} C_{pr\varrho}
\bigl[s_\tau^\varrho(stu)^{\mu+r-\varrho}\bigr]_{\widehat{uv}^{\lambda-r+\varrho}v^{\nu-\varrho},\,u^p,\,(vw=\widehat{xw})},
\]
and therefore does not permit the simplification that occurs already in the course of our calculation, once one knows that all terms with the factor $u_x$ are to be omitted.}:
\[
\tag{I}
s_x^m t_x^{n-\lambda}t_y^\lambda
=\sum_{\varrho}
\frac{\binom{m}{\varrho}\binom{\lambda}{\varrho}}
{\binom{m+n-\varrho+1}{\varrho}}
\bigl[(st\widehat{xy})^\varrho s_x^{m-\varrho}t_x^{n-\varrho}\bigr]_{y^{\lambda-\varrho}}.
\]

We combine the two expansions by setting up, for forms with two rows each of contragredient variables,
\[
s_x^m t_x^{n-\lambda}t_y^\lambda u_\sigma^\mu u_\tau^{\nu-\kappa}v_\tau^\kappa,
\]
expansions according to powers of $u_x$, whose coefficients are composite polars of the form sequence $s_x^m t_x^n u_\sigma^\mu u_\tau^\nu$, taken with respect to the variables $x$ and $u$.

It is enough to set up the series for two contragredient rows of symbols (respectively variables), since, by combining each two rows of symbols (variables) into a single new row, forms with more rows of symbols (variables) can be reduced to this case.

In order to ensure that all forms of the form sequence contain only two rows of symbols, respectively variables, we specialize:
\[
\begin{array}{ll}
\text{a) } \sigma=\widehat{st}, & \text{a') } y=\widehat{uv},\\[0.3em]
\text{b) } s=\widehat{\sigma\tau}, & \text{b') } v=\widehat{xy},
\end{array}
\]
and carry out the expansion for the specialization a), a').

By direct transfer of expansion I one obtains composite polars for forms $(stu)^\mu s_x^m t_x^{n-\lambda}t_y^\lambda$, whereas for forms $(stu)^\mu u_\tau^{\nu-\kappa}v_\tau^\kappa s_x^m t_x^{n-\lambda}t_y^\lambda$, instead of composite polars, terms of such a kind arise whose evaluation makes necessary the introduction of ever new auxiliary variables that are to be set equal only at the end, thereby complicating the calculation unnecessarily.

We therefore take an indirect route, by representing the composite polars of all forms of the form sequence, that is, expressions
\[
\bigl[s_\tau^\varrho(stu)^{\mu-\varrho+r}\bigr]_{y^\lambda}^{v^\kappa},
\]
by the initial members of these polars, and by inversion of the resulting recursive system of equations obtaining the desired expansion according to polars.

By the transfer principle, for the $\lambda$-th polar of a form $s_x^m t_x^n:[s_x^m t_x^n]_{y^\lambda}$ ($\lambda\le n$) the following expansion holds (the case $\lambda>n$ is reduced to the first by the expansion of $[s_y^m t_y^n]_{x^{m+n-\lambda}}$)\footnote{Gordan, \emph{Die Resultante binärer Formen} (Chap. I, § 5). Rend. Circ. Mat. di Palermo XXII. 1906.}:
\[
\bigl[s_x^m t_x^n\bigr]_{y^\lambda}
=\sum_r (-1)^r\,
\frac{\binom{m}{r}\binom{\lambda}{r}}{\binom{m+n}{r}}\,
(st\widehat{xy})^r s_x^{m-r}t_x^{n-\lambda}t_y^{\lambda-r},
\]
and correspondingly for the $\kappa$-th polar of $u_\sigma^\mu u_\tau^\nu:[u_\sigma^\mu u_\tau^\nu]_{v^\kappa}$
\[
\bigl[u_\sigma^\mu u_\tau^\nu\bigr]_{v^\kappa}
=\sum_\varrho (-1)^\varrho\,
\frac{\binom{\mu}{\varrho}\binom{\kappa}{\varrho}}{\binom{\mu+\nu}{\varrho}}\,
(\sigma\tau\widehat{uv})^\varrho u_\sigma^{\mu-\varrho}u_\tau^{\nu-\kappa}v_\tau^{\kappa-\varrho}.
\]

By multiplication it follows, with introduction of the specialization a) and a'),
\[
\begin{aligned}
\bigl[s_x^m t_x^n(stu)^\mu u_\tau^\nu\bigr]_{\widehat{uv}^{\lambda}}^{v^\kappa}
={}&
\sum_{r,\varrho}c_{r\varrho}\,
(st\widehat{uv})^r(st\widehat{uv})^\varrho
s_x^{m-r}t_x^{n-\lambda}(tuv)^{\lambda-r}(stu)^{\mu-\varrho}
u_\tau^{\nu-\kappa}v_\tau^{\kappa-\varrho},
\end{aligned}
\]
and from this, by expansion according to powers of $u_x$ with the first member distinguished ($\sum'_{r,\varrho}$ means that the member $r=0$, $\varrho=0$ is to be omitted) and taking account of (a) on p. 26,
\[
\tag{II}
\begin{aligned}
&s_x^m t_x^{n-\lambda}(tuv)^\lambda(stu)^\mu u_\tau^{\nu-\kappa}v_\tau^\kappa\\
&\quad =
\bigl[s_x^m t_x^n(stu)^\mu u_\tau^\nu\bigr]_{\widehat{uv}^{\lambda}}^{v^\kappa}\\
&\qquad
-\sum_{r,\varrho}' c_{r\varrho}\,s_\tau^\varrho(stu)^{\mu-\varrho+r}
u_\tau^{\nu-\kappa}v_\tau^{\kappa-\varrho}
s_x^{m-r}t_x^{n-\lambda}(tuv)^{\lambda-r+\varrho}v_x^r\\
&\qquad
+u_x\sum_{\varrho}\sum_{r=1}^{\lambda} c_{r\varrho}\,
s_\tau^\varrho(stu)^{\mu-\varrho+r-1}(stv)
u_\tau^{\nu-\kappa}v_\tau^{\kappa-\varrho}
s_x^{m-r}t_x^{n-\lambda}(tuv)^{\lambda-r+\varrho}v_x^{r-1}\\
&\qquad
\vdots\\
&\qquad
-(-1)^\lambda u_x^\lambda\sum_\varrho c_{\lambda\varrho}\,
s_\tau^\varrho(stu)^{\mu-\varrho}(stv)^\lambda
u_\tau^{\nu-\kappa}v_\tau^{\kappa-\varrho}
s_x^{m-\lambda}t_x^{n-\lambda}(tuv)^\varrho,
\end{aligned}
\]
where
\[
c_{r\varrho}=(-1)^{r+\varrho}\cdot
\frac{\binom{m}{r}\binom{\lambda}{r}}{\binom{m+n}{r}}\cdot
\frac{\binom{\mu}{\varrho}\binom{\kappa}{\varrho}}{\binom{\mu+\nu}{\varrho}}
\qquad
\begin{array}{l}
\text{for }\lambda\le n,\\
\kappa\le\nu.
\end{array}
\]
and correspondingly for the remaining combinations of the numbers $\lambda,n;\kappa,\nu$. Thus the expression $s_x^m t_x^{n-\lambda}(tuv)^\lambda(stu)^\mu u_\tau^{\nu-\kappa}v_\tau^\kappa$ is reduced to a composite polar, and, by applying the identity
\[
(stv)u_\tau=(stu)v_\tau-s_\tau(tuv),
\]
to a sum of analogously formed expressions corresponding to higher forms of the form sequence.

By iteration one arrives at the expansion
\[
\tag{III}
s_x^m t_x^{n-\lambda}(tuv)^\lambda(stu)^\mu u_\tau^{\nu-\kappa}v_\tau^\kappa
=
\sum_{p,r,\varrho}u_x^p\,C_{pr\varrho}\cdot
\bigl[s_\tau^\varrho(stu)^{\mu-\varrho+r}\bigr]_{\widehat{uv}^{\lambda-r+\varrho}v^{\kappa-\varrho+p},\,v_x^{r-p}}.
\]
The numerical coefficients $C_{pr\varrho}$ are computed uniquely and recursively from the coefficients $c_{r\varrho}$, $c'_{r\varrho}$ of the system of equations arising by iteration of (II) (cf. the example on p. 42). Analogous expansions arise for the other specializations.

\begin{center}
\textbf{§ 3. \quad Reduction theorems.}
\end{center}

The known theorems on the reduction of forms and form systems shall now be brought into connection with §§ 1 and 2, for which some definitions taken over from the theory of binary forms are necessary.

a) We call a system of forms a \emph{relatively complete system} modulo a prescribed sequence of forms if it has the property that all formations arising by folding an arbitrary product of these forms can be expressed as integral rational functions of the forms of the system, up to an additive term consisting of expressions that have arisen by folding the system forms with the system of the modulus\footnote{Gordan--Kerschensteiner, \emph{Vorlesungen über Invariantentheorie}. Vol. II. p. 227.}.

b) We call a form \emph{reducible} when it can be expressed by forms having invariants as factors or by ``higher forms,'' that is, by higher forms of the total form sequence to which the form belongs, or by forms containing the symbols of the modulus in higher order.

The theorems on reducers\footnote{Gordan, \emph{Math. Annalen} vol. XVII, p. 222.} can now be stated in their most general form as follows:

Definition: \emph{A reducer is a reducible form sequence.}

\emph{Theorem II. If the initial form of a form sequence is reducible because one of its members has arisen by folding with a reducer (has a reducer as factor), and if the final form of the form sequence has arisen from exactly this member by folding, then the total form sequence is reducible}\footnote{The simplification that results from Theorem II can be seen in the following example. For the special form $f=x_1^3x_2+x_2^3x_3+x_3^3x_1$, $a_y^2a_x^2u_y^2$ is a reducer. It follows by Theorem II that the form sequence
\[
\theta_\nu(\vartheta\nu x)\theta_x^3u_\vartheta u_\nu^2
=
a_\nu^2(abu)-a_\nu b_\nu(aba),
\]
is reduced, since $\theta_\nu^4=a_\nu^2b_\nu^2(abu)^2$ has arisen from the member $a_\nu^2(abu)$ by folding. In Gordan, loc. cit., p. 231, by contrast, the reduction of the forms
\[
\theta_\nu(\vartheta\nu x),\quad \theta_\nu(\vartheta\nu x)^2,\quad
\theta_\nu^2,\quad \theta_\nu^2(\vartheta\nu x),\quad
\theta_\nu^2(\vartheta\nu x)^2
\]
is carried out individually.}.

Proof: The form sequence arises from the initial member by folding in itself, hence by higher folding with the reducer on the one hand, and by folding the reducer in itself on the other hand. In both cases, by § 2, all forms of the form sequence can be represented as polars (transvections) over the form sequence of the reducer.

The inference from the initial form to the total form sequence no longer holds for the remaining reduction methods. These are:

1) so-called ``double reduction.''

By this we mean the reduction, in two different ways, of an expression that has two mutually independent reducers as factors, whereby relations arise between the higher forms to which one has reduced.

By systematic application of ``double reduction'' one must, on the one hand, obtain all relations\footnote{Thus, for example, the following were found by ``double reduction'': the reduction of the form $a_\eta^4$ (§ 10), the only form included superfluously in Maisano's system; and also the relations mentioned in the note to the Introduction between the three covariants and the three invariants.}; on the other hand, if ``double reduction'' applied to different expressions supplies no new relations, one has both a criterion for the independence of the higher forms and a check on the calculation.

2) \emph{Folding with decomposable forms.}

By ``decomposable forms'' are to be understood -- with a slight generalization of the usual concept -- forms that can be expressed by products of forms of lower degree and by ``higher'' forms. Products of forms pass into products under one folding; likewise products of nonlinear forms under twofold folding in the specializations a') and b') (§ 2), according to the product theorem:
\[
a_yb_y a_xb_x=\frac12a_y^2b_x^2+\frac12b_y^2a_x^2-\frac12(ab\widehat{yx})^2.
\]
First polars (transvections) and certain second polars of decomposable forms are therefore again decomposable forms. It follows that:

A member of a onefold (respectively twofold) transvection over a decomposable form, arising by onefold (respectively twofold) folding with a decomposable form, itself becomes a decomposable form:

a) if the next-higher forms in the form sequence of the decomposable form decompose or become reducible, or also if, under onefold folding, the specializations $y=\widehat{uv}$ or $v=\widehat{xy}$ occur;

b) if the remaining onefold foldings lead to higher forms (cf. § 12, B. $H_k$ and $H_k(k\eta x)$).

Such a member of a transvection can also decompose:

c) if the remaining onefold foldings lead to lower forms that decompose independently of the folding with the decomposable form, that is, that can be expressed by products and by the form to be reduced, although in each case one must first verify by calculation that the numerical coefficient of the member concerned does not vanish identically. This is nothing other than ``double reduction'' of the lower form (cf. § 23, C. $H_q(\theta H u)(\theta s u)$).

Decomposable forms arise by all onefold foldings with functional determinants\footnote{cf. Gordan, loc. cit. § 3.}, and moreover by certain higher foldings. One has
\[
(st\widehat{yx})s_y=\frac12\,t\,s_y^2-\frac12\,s\,t_y^2+\frac12(st\widehat{yx})^2,
\]
\[
(st\widehat{yx})^2s_y^2=\frac13\,t\,s_y^3-\frac13\,s\,t_y^3+s_y(st\widehat{yx})^2-\frac13(st\widehat{yx})^3.
\]
Analogous formulas hold for the dualistic forms
\[
(\sigma\tau\widehat{\nu u})v_\sigma
\quad\text{and}\quad
(\sigma\tau\widehat{\nu u})v_\sigma^2.
\]

From the theorems of this section the following rule follows for setting up all irreducible forms that arise from the folding of $S$ with $T$:

Beginning with the lowest foldings, proceed along any previously fixed path (for example row by row, or symmetrically with respect to the diagonal member) in the formation of the form sequence $ST$ until one reaches a form that has a reducer as factor. The form sequence defined by this form is to be omitted as soon as the final form satisfies the condition stated in Theorem II. After all reducible form sequences have been eliminated, all remaining reducible forms, and likewise all decomposable forms, are to be removed by applying double reduction. Places in the total form sequence that are covered twice by independently reducible form sequences give rise to the reduction of higher forms (cf. § 11, D).

\end{document}
